\subsection{\texorpdfstring{\(\mathfrak{sl}_2\) 三元组的定义}{\textbackslash mathfrak\{sl\}\_2 三元组的定义}}

定义 1。g 中的一个 \(\mathfrak{sl}_2\) 三元组是 g 中元素的一个序列 \(\mathfrak{g}\)\((x, h, y)\)\(\mathfrak{g}\)\((0, 0, 0)\)，它异于 ，并满足

\[
[ h, x ] = 2 x, [ h, y ] = - 2 y, [ x, y ] = - h.
\]

令 \((x,h,y)\) 为 g 中的一个 \(\mathfrak{sl}_{2}\) 三元组。从 \(\tau\) 到 g 的线性映射 \(\mathfrak{sl}(2,k)\)，满足 \(\tau(X_{+}) = x, \tau(H) = h, \tau(X_{-}) = y\)，是一个非零因而为单射的同态（因为 \(\mathfrak{sl}(2,k)\) 是单的），其像为 \(kx + kh + ky\)。于是我们得到从 g 中 \(\mathfrak{sl}_{2}\) 三元组的集合到从 \(\mathfrak{sl}(2,k)\) 到 g 的单射同态集合的一个典则双射。若 g 是半单的，且 \((x,h,y)\) 是 g 中的一个 \(\mathfrak{sl}_{2}\) 三元组，则 x、y 是 g 的幂零元素，而 h 是 g 的半单元素（第 I 章第6节第3号命题4）。

引理 1。令 \(x, h, y, y' \in \mathfrak{g}\)。若 \((x, h, y)\) 和 \((x, h, y')\) 是 \(\mathfrak{sl}_2\) 中的 \(\mathfrak{g}\) 三元组，则 \(y = y'\)。

事实上，\(y - y' \in \mathrm{Ker}(\mathrm{ad}_{\mathfrak{g}}x)\)，且 \((\mathrm{ad}_{\mathfrak{g}}h)(y - y') = -2(y - y')\)。但对于每个整数 \(\mathrm{ad}_{\mathfrak{g}}x\)，\(\mathrm{Ker}(p + \mathrm{ad}_{\mathfrak{g}}h)\) 在 \(p > 0\) 上都是单射（§1，第2号，命题2的推论）。

引理 2。令 \(\mathfrak{n}\) 为 \(\mathfrak{g}\) 的一个子代数，使得对每个 \(n\in \mathfrak{n}\)，\(\operatorname{ad}_{\mathfrak{g}}(n)\) 都是幂零的。令 \(h\in \mathfrak{g}\) 满足 \([h,\mathfrak{n}] = \mathfrak{n}\)。则 \(e^{\operatorname{ad}_{\mathfrak{g}}\mathfrak{n}}.h = h + \mathfrak{n}\)。

显然 \(e^{\operatorname{ad}_{\mathfrak{g}}(\mathfrak{n})}.h \subset h + \mathfrak{n}\)。我们将证明，若 \(v \in n\)，则 \(h + v \in e^{\operatorname{ad}_{\mathfrak{g}}(\mathfrak{n})}.h\)。只需证明对所有 \(h + v \in e^{\operatorname{ad}_{\mathfrak{g}}(\mathfrak{n})}.h + \mathscr {C} ^ {p} \mathfrak {n}\)，\(p \geq 1\)（因为当 p 足够大时 \(C^{p}n = 0\)）。当 p = 1 时这显然成立，因为 \(C^{1}n = n\)。现在假设我们已经证明存在 \(y_{p} \in n\) 和 \(z_{p} \in C^{p}n\)，使得 \(h + v = e^{\operatorname{ad}_{\mathfrak{g}}y_{p}}.h + z_{p}\)。由于 \((\operatorname{ad}_{\mathfrak{g}}h)(\mathfrak{n}) = \mathfrak{n}\)，\((\operatorname{ad}_{\mathfrak{g}}h)|\mathfrak{n}\) 是从 n 到 n 的双射；因此，它在保持 \(C^{p}n\) 不变的 \(C^{p}n\) 上的限制也是双射；于是存在 \(z \in C^{p}n\)，使得 \(z_{p} = [z, h]\)。于是

\[
e ^ {\operatorname{ad} _ {\mathfrak {g}} \left(y _ {p} + z\right)} h - e ^ {\operatorname{ad} _ {\mathfrak {g}} y _ {p}} h \in [ z, h ] + \mathscr {C} ^ {p + 1} \mathfrak {n}
\]

因此

\[
e ^ {\operatorname{ad} _ {\mathfrak {g}} \left(y _ {p} + z\right)} h \in h + v - z _ {p} + [ z, h ] + \mathscr {C} ^ {p + 1} \mathfrak {n} = h + v + \mathscr {C} ^ {p + 1} \mathfrak {n}
\]

这就完成了关于 p 的归纳证明。

引理 3。令 \(x \in \mathfrak{g}\)，\(\mathfrak{p} = \operatorname{Ker}(\operatorname{ad}x)\)，\(\mathfrak{q} = \operatorname{Im}(\operatorname{ad}x)\)。则 \([\mathfrak{p}, \mathfrak{q}] \subset \mathfrak{q}\)，且 \(\mathfrak{p} \cap \mathfrak{q}\) 是 \(\mathfrak{g}\) 的子代数。

若 \(u \in \mathfrak{p}\) 且 \(v \in \mathfrak{q}\)，则存在 \(w \in \mathfrak{g}\) 使 \(v = [x, w]\)，所以

\[
[ u, v ] = [ u, [ x, w ] ] = [ x, [ u, w ] ] - [ [ x, u ], w ] = [ x, [ u, w ] ] \in \mathfrak {q}.
\]

另一方面，p 是 g 的子代数，所以 \([p \cap q, p \cap q] \subset p \cap q\)。

引理 4。令 \((x,h,y)\) 和 \((x,h',y')\) 为 \(\mathfrak{sl}_2\) 中的 \(\mathfrak{g}\) 三元组。存在 \(z\in \mathfrak{g}\)，使得 \(\mathrm{ad}_{\mathfrak{g}}z\) 幂零，并且

\[
e ^ {\mathrm{ad} _ {\mathfrak {g}} z} x = x, \quad e ^ {\mathrm{ad} _ {\mathfrak {g}} z} h = h ^ {\prime}, \quad e ^ {\mathrm{ad} _ {\mathfrak {g}} z} y = y ^ {\prime}.
\]

令 \(\mathfrak{n} = \mathrm{Ker}(\mathrm{ad} x) \cap \mathrm{Im}(\mathrm{ad} x)\)。对所有 \(p \in \mathbf{Z}\)，令 \(\mathfrak{g}_p = \mathrm{Ker}(\mathrm{ad} h - p)\)。由 §1，第3号（应用于 \(kx + ky + kh\) 在 \(\mathfrak{g}\) 上的伴随表示），有 \(\mathfrak{n} \subset \sum_{p > 0} \mathfrak{g}_p\)，所以对所有 \(\mathrm{ad}_{\mathfrak{g}} n\)，\(n \in \mathfrak{n}\) 都是幂零的，且 \([h, \mathfrak{n}] = \mathfrak{n}\)。我们有 \([x, h' - h] = 0\) 以及 \([x, y - y'] = h' - h\)，所以 \(h' - h \in \mathfrak{n}\)。由引理2和3，存在 \(z \in \mathfrak{n}\)，使 \(e^{\mathrm{ad}_{\mathfrak{g}} z} h = h'\)。由于 \(z \in \mathrm{Ker} \mathrm{ad}_{\mathfrak{g}} x\)，有 \(e^{\mathrm{ad}_{\mathfrak{g}} z} x = x\)。现在由引理1 得 \(e^{\mathrm{ad}_{\mathfrak{g}} z} y = y'\)。证毕。

令 G 为 g 的自同构群。若存在  使 、、，则称 g 中的两个 \(sl_{2}\) 三元组 \((x,h,y)\)、\((x',h',y')\) G-共轭。\(g\in G\)\(gx=x'\)\(gh=h',gy=y'\)

命题 1。令 G 为包含 \(\mathrm{Aut}_e(\mathfrak{g})\) 的 g 的自同构群。令 \((x,h,y)\) 和 \((x',h',y')\) 为 g 中的 \(\mathfrak{sl}_2\) 三元组。令

\[
\mathfrak {t} = k x + k h + k y, \quad \mathfrak {t} ^ {\prime} = k x ^ {\prime} + k h ^ {\prime} + k y ^ {\prime}.
\]

考虑以下条件：

\begin{enumerate}
\def\labelenumi{(\roman{enumi})}
\item
  x 与 \(x'\) G-共轭；
\item
  \((x, h, y)\) 与 \((x', h', y')\) G-共轭；
\item
  \(\mathfrak{t}\) 与 \(\mathfrak{t}'\) G-共轭。
\end{enumerate}

有 (i) \(\Longleftrightarrow\) (ii) \(\Longrightarrow\) (iii)。若 \(k\) 是代数闭域，则这三个条件等价。

\begin{enumerate}
\def\labelenumi{(\roman{enumi})}
\item
  \(\Longleftrightarrow\) (ii)：这由引理4 得出。
\item
  \(\Longrightarrow\) (iii)：显然。
\end{enumerate}

我们假设 k 是代数闭域，并证明 (iii) \(k\)\(\Longrightarrow\) (i)。先处理 \(\mathfrak{t} = \mathfrak{t}' = \mathfrak{g} = \mathfrak{sl}(2, k)\) 的情形。由于 \(\operatorname{ad}_{\mathfrak{g}} x\) 是幂零的，k² 的自同态 x 是幂零的（第 I 章第6节定理3），所以存在矩阵 \(x\)\(k^2\)\(A \in \mathbf{GL}(2, k)\)，使 \(AxA^{-1} = X\)，从而存在  的自同构 \(\alpha\)，使 \(\mathfrak{sl}(2, k)\)\(\alpha(x) = x'\)；现在 \(\alpha \in \operatorname{Aut}_e(\mathfrak{g})\)（§5，第3号，命题5的推论2）。现在转到一般情形；假设 \(\mathfrak{t}\) 和 \(\mathfrak{t}'\) G-共轭，并证明 x 和 \(x\)\(x'\) G-共轭。可以假设 \(\mathfrak{t} = \mathfrak{t}'\)。由上文，存在 \(\beta \in \operatorname{Aut}_e(\mathfrak{t})\)，使 \(\beta x = x'\)。现在，若 \(t \in \mathfrak{t}\) 满足 \(\operatorname{ad}_{\mathfrak{t}} t\) 幂零，则 \(\operatorname{ad}_{\mathfrak{g}} t\) 也幂零；所以 \(\beta\) 延拓为 \(\operatorname{Aut}_e(\mathfrak{g})\) 的一个元素。

注。只假设 \(k = k^2\) 时，命题1 的三个条件仍然等价（参见~习题1）。

